% a01_H0rel part 6x3: the replaced step (d) and the parameter conditions (input from a01_H0rel_p6.tex after a01_p6x2)
\subparagraph{(d) The partition \cite[(3.43)--(3.45)]{Ba2}.} The printed ``regular partition into a lattice of cubes, each a sum of large blocks, with side in $(r,2r)$'' needs, per coordinate, a divisor $c$ of the number $N$ of large blocks with $r/s_b<c<2r/s_b$; at $L=2$ a power of $2$ always works, while at $L=5$ the divisors of $N_i=2MJ/(\ell_is_b)$, a number of the form $2\cdot5^{b}J$, have gaps of ratio up to $\frac52$, so the window can contain no divisor, and with integer sides and even $N$ the printed partition fails exactly when $r_i/s_b$ is an integer. The partition is used in \cite[Sect.~3C]{Ba2} only through four properties: the covering \cite[(3.43)]{Ba2}; the neighbour count $3^d$ in \cite[(3.44)--(3.45)]{Ba2}; the volume bounds $(2r)^d$ in \cite[(3.48)--(3.49)]{Ba2}; and the corridor step of \cite[p.~593]{Ba2}.

\begin{lemma}[the product partition and the cube count]\label{lem:a01:PP}
Let $q_i\defeq\lceil(r_i+s_b)/s_b\rceil$ and let every coordinate carry $N\ge q_i$ large blocks (on King's cube, $N=2MJ/(\ell_is_b)$ large blocks per side in scale-$i$ units, which exceeds $q_i$ for $J\ge J_0$; on a torus, $N$ per period). Cut each coordinate interval into consecutive arcs of between $q_i$ and $2q_i-1$ large blocks (write $N=nq_i+\varrho'$, $0\le\varrho'<q_i$, and give the $\varrho'$ extra blocks to the $n$ arcs, at most $q_i-1$ each), and let $\mathcal Q_i$ be the product partition into boxes. Call $Q\sim Q'$ if their arc indices differ by at most one in every coordinate. Then: \textup{(Q1)} every box is a union of large blocks; \textup{(Q2)} every box is a product of intervals of at most $2r_i+3s_b\le(2+3s_b/R)r_i$ points; \textup{(Q3a)} every box has at most $3^d$ boxes $Q'\sim Q$; \textup{(Q3b)} if $\dist(z,x)\le r_i+s_b$ then $Q(z)\sim Q(x)$. Consequently, with $I_i\defeq\{z:\dist(z,(\Lambda_7^{(i-1)\prime})^c)\ge s_b\}$ and $\mathcal C_i\defeq\{Q\in\mathcal Q_i:Q\cap\Lambda_0^{(i)c}\cap I_i\ne\emptyset\}$, every $Q\in\mathcal C_i$ satisfies $Q\sim Q(x_e)$ for the charge point $x_e$ of some element $e$ of the family $F_i$ of large-field events at scale $i$, hence $|\mathcal C_i|\le3^d|F_i|$; the volume bound $(2r)^d$ becomes $(2r_i+3s_b)^d$ and the corridor bound $22r$ becomes $22r_i+3s_b$, both $O(1)r_i^d$.
\end{lemma}
\begin{proof}
(Q1) Arcs are unions of large-block intervals. (Q2) An arc has at most $(2q_i-1)s_b<2(r_i+s_b)+s_b$ points, and $s_b\le(s_b/R)r_i$. (Q3a) The boxes $\sim Q$ are the $3^d$ neighbours of $Q$ in the product structure. (Q3b) Every arc has $\ge q_is_b\ge r_i+s_b$ points, so if $|z-x|_\infty\le r_i+s_b$ the two arcs in each coordinate are equal or consecutive (otherwise a whole arc lies strictly between and the coordinate distance is $\ge q_is_b+1$). For the count: let $Q\in\mathcal C_i$ and $z\in Q\cap\Lambda_0^{(i)c}\cap I_i$; its large block $\beta$ lies in $Q$ by (Q1) and, having sup-diameter $<s_b$, in $\Lambda_7^{(i-1)\prime}$; since $z\notin\Lambda_0^{(i)}$, $\beta$ is not one of the blocks of \cite[p.~570]{Ba2}, so $\dist(\beta,S_i)\le r_i$ where $S_i$ is the point set of $F_i$; take $z'\in\beta$, $y\in S_i$ with $\dist(z',y)\le r_i$, $y$ a point of an element $e$ with $\dist(y,x_e)\le L-1<s_b$; then $\dist(z',x_e)\le r_i+s_b$ and $Q=Q(z')\sim Q(x_e)$ by (Q3b); count with (Q3a). The printed \cite[(3.45)]{Ba2} and the first inequality of \cite[(3.46)]{Ba2} do not hold as printed (a large block straddling $\partial\Lambda_7^{(i-1)\prime}$ lies in $\Lambda_0^{(i)c}$ with $F_i$ possibly empty); the count above, on the interior $I_i$, becomes the count of D14 over the paid elements through \cref{lem:a01:PP5}; in D14 the boxes of $\mathcal C_{i'}$ for $i'<i$ are carried along by the corridor construction of \cite[(3.48)--(3.50)]{Ba2} with the thickness $7(r_{i-1}+s_b)+Ls_b$ inside Ba{\l}aban's $(9r,10r)$ once $(7+L)s_b<2r_{i-1}$.
\end{proof}
% a01_H0rel part 6x3, continued: the count over the paid elements (input from a01_p6x3.tex after lem:a01:PP)
\begin{lemma}[the count over the paid elements]\label{lem:a01:PP5}
Let the boxes $\mathcal Q_i$, the relation $\sim$, the arcs of $q_i$ to $2q_i-1$ large blocks, the set $\mathcal C_i$ and the family $F_i$ be those of \cref{lem:a01:PP}, with the charge point $x_e$ of an element $e$ as in its proof (every point of $e$ lies within $\dist\le L-1$ of $x_e$). Let $F_i^{\mathrm{paid}}\subset F_i$, and suppose that every $e\in F_i\setminus F_i^{\mathrm{paid}}$ has a \emph{witness} $w(e)\in F_i^{\mathrm{paid}}$ with $\dist(x_e,x_{w(e)})\le r_i+s_b$. Then
\[
|\mathcal C_i|\ \le\ 5^d\,|F_i^{\mathrm{paid}}|.
\]
In D14 the unpaid elements are the near $R_s$ elements, the blocks $y\in R_s^{(i)}$ within $\ell^\infty$-distance $R_0$ of the face layer, and $F_i^{\mathrm{paid}}$ consists of Ba{\l}aban's six families, $R_s$ restricted to its far elements, and the face family; the witness of D7(b) (\cref{lem:a01:F541}(e)) has a point within distance $dR_0+1$ of a point of $e$, so $\dist(x_e,x_{w(e)})\le dR_0+1+2(L-1)\le s_b+L-2<r_i+s_b$ by $s_b\ge dR_0+L+1$ (D7(c)) and $r_i\ge R>L$, and the hypothesis holds.
\end{lemma}
\begin{proof}
Let $Q\in\mathcal C_i$. The proof of \cref{lem:a01:PP} gives $z'\in Q$ and $e\in F_i$ with $\dist(z',x_e)\le r_i+s_b$. Put $e'\defeq e$ if $e\in F_i^{\mathrm{paid}}$ and $e'\defeq w(e)$ otherwise; in both cases $e'\in F_i^{\mathrm{paid}}$ and $\dist(z',x_{e'})\le2(r_i+s_b)\le2q_is_b$, since $q_is_b\ge r_i+s_b$. In each coordinate the arc indices of $Q=Q(z')$ and of $Q(x_{e'})$ differ by at most $2$: if they differed by $3$ or more, two whole arcs, of at least $q_is_b$ points each, would lie strictly between the two coordinates, whose distance would then be $\ge2q_is_b+1$. A box has at most $5^d$ boxes whose arc indices differ from its own by at most $2$ in every coordinate (fewer on a torus with fewer than five arcs in some coordinate). So $Q\mapsto e'$ maps $\mathcal C_i$ into $F_i^{\mathrm{paid}}$ with fibres of at most $5^d$ boxes, which is the claim. For the application: a near $R_s$ element $e$ at the block $y$ has $\dist(y,x_e)\le L-1$; \cref{lem:a01:F541}(e) gives, on a block path of length $\le dR_0$ from $y$ to the face layer, a bond of the gradient family or a block of the face family, that is an element $w(e)$ of a paid family one of whose points lies within distance $dR_0+1$ of $y$ and within $L-1$ of $x_{w(e)}$; the triangle inequality gives $\dist(x_e,x_{w(e)})\le dR_0+2L-1$, and $dR_0+2L-1\le s_b+L-2$ is $s_b\ge dR_0+L+1$.
\end{proof}


The replaced objects of (a)--(d), $r^\flat$, $D^\flat$, $\zeta^{(i)}_\theta$, the box $\square$ of \cref{lem:a01:lem24} and the partition $\mathcal Q_i$, are built from $(L,M,R,\ell_i,\bar n)$ and the set of event blocks alone, so they are the same on the cube as on the torus (\cref{thm:a06:BALA}); they are auxiliary objects of the proofs, not of the statements, and under (b) the all-small-field action changes by $O(\ell_i^\kappa)|X^{(i)}|$, which the statements absorb. The parameter conditions they add are lower bounds on $R$ ($R\ge R^\flat$; $R\ge10(M+1)$, $R\ge M_\theta+2$; $R\ge10s_b+10$; $R\ge6L\sqrt d$ for the covering of D14) and caps on $\eps_0^{\mathrm{Ba}}$, all fixed in Ba{\l}aban's order of choices ($d$, $a$, $\bar n$, $\alpha_0$, $L$, $M$, $R$, $\varrho$, $p$, $b_0$, then $\eps_0^{\mathrm{Ba}}$), so the order is acyclic.
